\DocumentMetadata{
  pdfversion=2.0,
  pdfstandard=ua-2,
  lang=en-US,
  tagging=on,
  tagging-setup={math/setup=mathml-SE,tabsorder=structure}
}
\documentclass[11pt,letterpaper]{article}
\usepackage[margin=0.68in,includefoot]{geometry}
\usepackage{newcomputermodern}
\usepackage{amsmath,amscd,mathtools}
\usepackage{tikz,adjustbox}
\providecommand{\square}{\mdlgwhtsquare}
\usepackage[hidelinks]{hyperref}
\hypersetup{
  pdftitle={Analysis First-Year Exam, August 2019, Part 1},
  pdfauthor={University of Florida Department of Mathematics},
  pdfsubject={Graduate mathematics examination},
  pdfdisplaydoctitle=true,
  bookmarksopen=true
}
\setcounter{secnumdepth}{0}
\setlength{\parindent}{0pt}
\setlength{\parskip}{0.28em}
\setlength{\emergencystretch}{3em}
\setlength{\leftmargini}{2.15em}
\setlength{\leftmarginii}{2.65em}
\setlength{\leftmarginiii}{3em}
\pagestyle{empty}
\raggedbottom
\allowdisplaybreaks
\definecolor{ProblemConcern}{HTML}{7A1F1F}
\newenvironment{problemconcern}{%
  \par\tagstructbegin{tag=Aside}\begingroup\color{ProblemConcern}
}{%
  \par\endgroup\tagstructend\nopagebreak[4]\vspace{0.12em}
}
\NewTaggingSocketPlug{block/list/label}{examproblem}{%
  \tagstructbegin{tag=Lbl}\tagstructend%
  \tagstructbegin{tag=\LBody}%
  \tagstructbegin{tag=\ProblemHeadingTag,title-o={\ProblemHeadingText}}%
  \tagmcbegin{tag=\ProblemHeadingTag,actualtext={\ProblemHeadingText}}%
  #2%
  \tagmcend%
  \tagstructend%
}
\newenvironment{examproblems}[2]{%
  \def\ProblemHeadingTag{#2}%
  \AssignTaggingSocketPlug{block/list/label}{examproblem}%
  \begin{enumerate}%
  \setcounter{enumi}{#1}%
  \renewcommand{\labelenumi}{\textbf{\theenumi.}}%
  \setlength{\topsep}{0.35em}%
  \setlength{\partopsep}{0pt}%
  \setlength{\itemsep}{0.48em}%
  \setlength{\parsep}{0.18em}%
}{\end{enumerate}}
\newenvironment{examsubparts}{%
  \AssignTaggingSocketPlug{block/list/label}{default}%
  \begin{enumerate}%
  \setlength{\topsep}{0.18em}%
  \setlength{\partopsep}{0pt}%
  \setlength{\itemsep}{0.16em}%
  \setlength{\parsep}{0.1em}%
}{\end{enumerate}}
\newenvironment{examinstructions}{%
  \par\begingroup\itshape
}{%
  \par\endgroup\vspace{0.18em}
}
\makeatletter
\renewcommand\subsection{\@startsection{subsection}{2}{0pt}%
  {-1.25ex plus -.25ex minus -.1ex}%
  {0.55ex plus .1ex}%
  {\normalfont\large\bfseries}}
\renewcommand{\@maketitle}{%
  \newpage\null\vskip -1em%
  {\footnotesize\bfseries
    \href{https://www.ufl.edu/}{University of Florida}, \href{https://math.ufl.edu/}{Department of Mathematics}\par}%
  \vskip -0.5em%
  \tagstructbegin{tag=H1,title={Analysis First-Year Exam, August 2019, Part 1}}%
  \tagpdfparaOff%
  \tagmcbegin{tag=H1}%
    {\LARGE\bfseries\href{https://gma.math.ufl.edu/exams/analysis-first-year/}{Analysis First-Year Exam}, August 2019, Part 1\par}%
  \tagmcend%
  \tagpdfparaOn%
  \tagstructend%
  \vskip 0.25em%
  \tagmcbegin{artifact}%
    \hrule height 1pt%
  \tagmcend%
  \vskip 1em%
}
\makeatother
\title{Analysis First-Year Exam, August 2019, Part 1}
\author{}
\date{}
\begin{document}
\maketitle
\thispagestyle{empty}
\begin{examinstructions}
Do exactly 4 problems. Work must be presented in a neat and logical fashion in order to receive credit. Do not leave any gaps. When a theorem is used in a proof it must be precisely stated.
\end{examinstructions}
\begin{examproblems}{0}{H2}
\def\ProblemHeadingText{Problem 1.}
\item
This problem has two parts.
\begin{examsubparts}
\item[\textnormal{a)}]
Let~\(X\) be a metric space, \((p_n)\) a sequence in~\(X\), and \(p\in X\). Prove that \((p_n)\) converges to~\(p\) if and only if every subsequence of \((p_n)\) has a subsequence converging to~\(p\).
\item[\textnormal{b)}]
Give an example of a metric space~\(X\) and a sequence \((p_n)\) in~\(X\) such that \((p_n)\) diverges, but every subsequence of \((p_n)\) has a convergent subsequence.
\end{examsubparts}
\def\ProblemHeadingText{Problem 2.}
\item
Let \(f:(0,1)\to\mathbb{R}\) be a uniformly continuous function. Prove that \(\lim_{x\to 0^+}f(x)\) exists. Does the conclusion remain true if~\(f\) is only assumed bounded and continuous?
\def\ProblemHeadingText{Problem 3.}
\item
Let~\(X\) be a metric space and~\(S\) a connected subset of~\(X\). Must \(\overline{S}\) (the closure of~\(S\)) be connected? Prove, or give a counterexample. If \(S^\circ\) (the interior of~\(S\)) is nonempty, must it be connected? Prove, or give a counterexample.
\def\ProblemHeadingText{Problem 4.}
\item
Let~\(X\) and~\(Y\) be disjoint, nonempty subsets of \(\mathbb{R}^n\), with~\(X\) closed and~\(Y\) compact. Define
\[
\operatorname{dist}(X,Y):=\inf_{x\in X,\,y\in Y}\lVert x-y\rVert.
\]
 Prove that there exist points \(x_0\in X\), \(y_0\in Y\) with \(\lVert x_0-y_0\rVert=\operatorname{dist}(X,Y)\).
\def\ProblemHeadingText{Problem 5.}
\item
\begin{problemconcern}
\textbf{Warning: Suspected error.} The problem is preserved as printed, but its conclusion generally fails without an additional condition such as \(f(0)=0\), or without replacing \(f(x)/x\) by \((f(x)-f(0))/x\). The intended repair is not uniquely determined.
\end{problemconcern}
Let \(f:[0,a)\to\mathbb{R}\) be differentiable, and suppose that \(f'\) is strictly increasing on \([0,a)\). Prove that the function
\[
g(x)=
\begin{cases}
f'(0), & x=0,\\
\dfrac{f(x)}{x}, & 0<x<a
\end{cases}
\]
 is continuous and strictly increasing on \([0,a)\).
\end{examproblems}
\end{document}
